% GENERATED FILE. Edit canonical lessons, then run npm run generate:books.
% canonical-source-hash: fnv1a64:f53a396a104dbbbe
% canonical-lessons: ML-C384-jananam, ML-C384-maranam, ML-C384-caramam, ML-C384-atmakatha, ML-C384-jivacaritram

\chapter{Births, Deaths and Life Stories}
\label{ch:ml-births-deaths-and-life-stories}
\hlchaptermodality{384}

\begin{chapteropening}
\textbf{By the end of this chapter:} \emph{I can say birth, death, an obituary, an autobiography and a biography.}

Complete the last lesson of chapter 384: I can say birth, death, an obituary, an autobiography and a biography.
\end{chapteropening}

\section[\texorpdfstring{\ml{ജനനം} (jananaṁ)}{ജനനം (jananaṁ)}]{\ml{ജനനം} (jananaṁ) — birth}

\label{lesson:ML-C384-jananam}



\begin{quote}
Before the new one: say the Malayalam for to find out, then the Malayalam for to retrieve.
\end{quote}

\subsection*{You'll want to know: \ml{ജനനം}}
\textbf{\ml{ജനനം}} — \emph{jananaṁ} — \textquotedblleft{}birth\textquotedblright{}.

\begin{quote}
\ml{കുഞ്ഞിന്റെ} \ml{ജനനം} \ml{വീട്ടിൽ} \ml{എല്ലാവർക്കും} \ml{വലിയ} \ml{സന്തോഷമായി}.
\end{quote}

\begin{quote}
\emph{kuññinṟe jananaṁ vīṭṭil ellāvarkkuṁ valiya santōṣamāyi.}
\end{quote}

\begin{quote}
The baby's birth was a great joy for everyone at home.
\end{quote}

The verb is \ml{ജനിക്കുക}, to be born.

\begin{cousinweb}
\textbf{\ml{ജനനം}} (\emph{jananaṁ}) is borrowed from Sanskrit \emph{janana}, \textquotedblleft{}birth\textquotedblright{} (CDIAL 5099), from the root \emph{jan}, \textquotedblleft{}to be born\textquotedblright{}, the root of \ml{ജനിക്കുക} (\emph{janikkuka}).
\end{cousinweb}

\begin{grammarlens}[title={What you've built}]
One more word for telling what happened.
\end{grammarlens}

\subsection*{Your turn}
\begin{itemize}
\raggedright
  \item \emph{Say it:} \emph{jananaṁ}
  \item \emph{Say it:} \emph{jananaṁ}, once more
  \item \emph{Say it:} \emph{kuññinṟe jananaṁ vīṭṭil ellāvarkkuṁ valiya santōṣamāyi.} — The baby's birth was a great joy for everyone at home.
  \item \emph{Recall:} say the Malayalam for to find out, then the Malayalam for to retrieve, then say \emph{jananaṁ} again
\end{itemize}

\begin{tcolorbox}[breakable,colback=teal!4,colframe=teal!35!black,title={Before you move on}]
What does \ml{ജനനം} mean? (\textquotedblleft{}Birth\textquotedblright{}.) Say it once more.
\end{tcolorbox}
\section[\texorpdfstring{\ml{മരണം} (maraṇaṁ)}{മരണം (maraṇaṁ)}]{\ml{മരണം} (maraṇaṁ) — death}

\label{lesson:ML-C384-maranam}



\begin{quote}
Before the new one: say the Malayalam for to retrieve, then the Malayalam for birth.
\end{quote}

\subsection*{You'll want to know: \ml{മരണം}}
\textbf{\ml{മരണം}} — \emph{maraṇaṁ} — \textquotedblleft{}death\textquotedblright{}.

\begin{quote}
\ml{മുത്തച്ഛന്റെ} \ml{മരണം} \ml{എല്ലാവരെയും} \ml{സങ്കടപ്പെടുത്തി}.
\end{quote}

\begin{quote}
\emph{muttacchanṟe maraṇaṁ ellāvareyuṁ saṅkaṭappeṭutti.}
\end{quote}

\begin{quote}
Grandfather's death saddened everyone.
\end{quote}

The verb is \ml{മരിക്കുക}, to die.

\begin{cousinweb}
\textbf{\ml{മരണം}} (\emph{maraṇaṁ}) is borrowed from Sanskrit \emph{maraṇa}, \textquotedblleft{}dying, death\textquotedblright{} (CDIAL 9869), the noun beside \ml{മരിക്കുക} (\emph{marikkuka}), \textquotedblleft{}to die\textquotedblright{}.
\end{cousinweb}

\begin{grammarlens}[title={What you've built}]
One more word for telling what happened.
\end{grammarlens}

\subsection*{Your turn}
\begin{itemize}
\raggedright
  \item \emph{Say it:} \emph{maraṇaṁ}
  \item \emph{Say it:} \emph{maraṇaṁ}, once more
  \item \emph{Say it:} \emph{muttacchanṟe maraṇaṁ ellāvareyuṁ saṅkaṭappeṭutti.} — Grandfather's death saddened everyone.
  \item \emph{Recall:} say the Malayalam for to retrieve, then the Malayalam for birth, then say \emph{maraṇaṁ} again
\end{itemize}

\begin{tcolorbox}[breakable,colback=teal!4,colframe=teal!35!black,title={Before you move on}]
What does \ml{മരണം} mean? (\textquotedblleft{}Death\textquotedblright{}.) Say it once more.
\end{tcolorbox}
\section[\texorpdfstring{\ml{ചരമം} (caramaṁ)}{ചരമം (caramaṁ)}]{\ml{ചരമം} (caramaṁ) — an obituary; a death}

\label{lesson:ML-C384-caramam}



\begin{quote}
Before the new one: say the Malayalam for birth, then the Malayalam for death.
\end{quote}

\subsection*{You'll want to know: \ml{ചരമം}}
\textbf{\ml{ചരമം}} — \emph{caramaṁ} — \textquotedblleft{}an obituary; a death\textquotedblright{}.

\begin{quote}
\ml{അച്ഛൻ} \ml{എന്നും} \ml{രാവിലെ} \ml{പത്രത്തിലെ} \ml{ചരമം} \ml{വായിക്കും}.
\end{quote}

\begin{quote}
\emph{acchan ennuṁ rāvile patrattile caramaṁ vāyikkuṁ.}
\end{quote}

\begin{quote}
Every morning father reads the obituaries in the paper.
\end{quote}

Literally the last: a formal word for a death, and for the death notices in a newspaper.

\begin{cousinweb}
\textbf{\ml{ചരമം}} (\emph{caramaṁ}) is borrowed from Sanskrit \emph{carama}, \textquotedblleft{}last, ultimate, final\textquotedblright{} (Monier-Williams): the last event of a life.
\end{cousinweb}

\begin{grammarlens}[title={What you've built}]
One more word for telling what happened.
\end{grammarlens}

\subsection*{Your turn}
\begin{itemize}
\raggedright
  \item \emph{Say it:} \emph{caramaṁ}
  \item \emph{Say it:} \emph{caramaṁ}, once more
  \item \emph{Say it:} \emph{acchan ennuṁ rāvile patrattile caramaṁ vāyikkuṁ.} — Every morning father reads the obituaries in the paper.
  \item \emph{Recall:} say the Malayalam for birth, then the Malayalam for death, then say \emph{caramaṁ} again
\end{itemize}

\begin{tcolorbox}[breakable,colback=teal!4,colframe=teal!35!black,title={Before you move on}]
What does \ml{ചരമം} mean? (\textquotedblleft{}An obituary; a death\textquotedblright{}.) Say it once more.
\end{tcolorbox}
\section[\texorpdfstring{\ml{ആത്മകഥ} (ātmakatha)}{ആത്മകഥ (ātmakatha)}]{\ml{ആത്മകഥ} (ātmakatha) — an autobiography}

\label{lesson:ML-C384-atmakatha}



\begin{quote}
Before the new one: say the Malayalam for death, then the Malayalam for an obituary.
\end{quote}

\subsection*{You'll want to know: \ml{ആത്മകഥ}}
\textbf{\ml{ആത്മകഥ}} — \emph{ātmakatha} — \textquotedblleft{}an autobiography\textquotedblright{}.

\begin{quote}
\ml{ആ} \ml{നടന്റെ} \ml{ആത്മകഥ} \ml{ഒരുപാട്} \ml{പേർ} \ml{വായിച്ചു}.
\end{quote}

\begin{quote}
\emph{ā naṭanṟe ātmakatha orupāṭŭ pēr vāyiccu.}
\end{quote}

\begin{quote}
Many people read that actor's autobiography.
\end{quote}

\ml{ആത്മ}-, self, + \ml{കഥ}, a story: the story of oneself.

\begin{cousinweb}
\textbf{\ml{ആത്മകഥ}} (\emph{ātmakatha}) is built from two Sanskrit words, \emph{ātman}, \textquotedblleft{}breath, soul, self\textquotedblright{} (CDIAL 1135), and \emph{kathā}, \textquotedblleft{}speech, a story\textquotedblright{} (CDIAL 2704), the source of \ml{കഥ} (\emph{katha}).
\end{cousinweb}

\begin{grammarlens}[title={What you've built}]
One more word for telling what happened.
\end{grammarlens}

\subsection*{Your turn}
\begin{itemize}
\raggedright
  \item \emph{Say it:} \emph{ātmakatha}
  \item \emph{Say it:} \emph{ātmakatha}, once more
  \item \emph{Say it:} \emph{ā naṭanṟe ātmakatha orupāṭŭ pēr vāyiccu.} — Many people read that actor's autobiography.
  \item \emph{Recall:} say the Malayalam for death, then the Malayalam for an obituary, then say \emph{ātmakatha} again
\end{itemize}

\begin{tcolorbox}[breakable,colback=teal!4,colframe=teal!35!black,title={Before you move on}]
What does \ml{ആത്മകഥ} mean? (\textquotedblleft{}An autobiography\textquotedblright{}.) Say it once more.
\end{tcolorbox}
\section[\texorpdfstring{\ml{ജീവചരിത്രം} (jīvacaritraṁ)}{ജീവചരിത്രം (jīvacaritraṁ)}]{\ml{ജീവചരിത്രം} (jīvacaritraṁ) — a biography}

\label{lesson:ML-C384-jivacaritram}



\begin{quote}
Before the new one: say the Malayalam for an obituary, then the Malayalam for an autobiography.
\end{quote}

\subsection*{You'll want to know: \ml{ജീവചരിത്രം}}
\textbf{\ml{ജീവചരിത്രം}} — \emph{jīvacaritraṁ} — \textquotedblleft{}a biography\textquotedblright{}.

\begin{quote}
\ml{ഞാൻ} \ml{ഒരു} \ml{കവിയുടെ} \ml{ജീവചരിത്രം} \ml{വായിക്കുന്നു}.
\end{quote}

\begin{quote}
\emph{ñān oru kaviyuṭe jīvacaritraṁ vāyikkunnu.}
\end{quote}

\begin{quote}
I am reading a poet's biography.
\end{quote}

\ml{ജീവ}-, life, + \ml{ചരിത്രം}, history: the history of one life.

\begin{cousinweb}
\textbf{\ml{ജീവചരിത്രം}} (\emph{jīvacaritraṁ}) is built from Sanskrit \emph{jīva}, \textquotedblleft{}life\textquotedblright{} (CDIAL 5239), and \emph{caritra}, \textquotedblleft{}going; behaviour\textquotedblright{} (CDIAL 4691), the source of \ml{ചരിത്രം} (\emph{caritraṁ}), \textquotedblleft{}history\textquotedblright{}.
\end{cousinweb}

\begin{grammarlens}[title={What you've built}]
One more word for telling what happened.
\end{grammarlens}

\subsection*{Your turn}
\begin{itemize}
\raggedright
  \item \emph{Say it:} \emph{jīvacaritraṁ}
  \item \emph{Say it:} \emph{jīvacaritraṁ}, once more
  \item \emph{Say it:} \emph{ñān oru kaviyuṭe jīvacaritraṁ vāyikkunnu.} — I am reading a poet's biography.
  \item \emph{Recall:} say the Malayalam for an obituary, then the Malayalam for an autobiography, then say \emph{jīvacaritraṁ} again
\end{itemize}

\begin{tcolorbox}[breakable,colback=teal!4,colframe=teal!35!black,title={Before you move on}]
What does \ml{ജീവചരിത്രം} mean? (\textquotedblleft{}A biography\textquotedblright{}.) Say it once more.
\end{tcolorbox}
